\documentclass[UTF8,11pt]{ctexart}
\usepackage[a4paper,margin=25mm]{geometry}
\usepackage{amsmath,amssymb,amsthm,mathtools}
\usepackage[all]{xy}
\usepackage{xurl}
\usepackage{hyperref}
\usepackage{enumitem}
\hypersetup{hidelinks,breaklinks=true}
\setlength{\emergencystretch}{3em}
\newtheorem*{theorem}{Theorem}
\newtheorem*{lemma}{Lemma}
\newtheorem*{proposition}{Proposition}
\newtheorem*{corollary}{Corollary}
\newtheorem*{definition}{Definition}
\newcommand{\Spec}{\operatorname{Spec}}
\newcommand{\Hom}{\operatorname{Hom}}
\newcommand{\End}{\operatorname{End}}
\newcommand{\Gal}{\operatorname{Gal}}
\newcommand{\Cl}{\operatorname{Cl}}
\newcommand{\vol}{\operatorname{vol}}
\newcommand{\covol}{\operatorname{covol}}
\newcommand{\tr}{\operatorname{tr}}
\newcommand{\im}{\operatorname{im}}
\newcommand{\id}{\operatorname{id}}
\newcommand{\coker}{\operatorname{coker}}
\newcommand{\stacksref}[2]{\href{https://stacks.math.columbia.edu/tag/#1}{#2 [#1]}}
\begin{document}
\section*{002\quad 理想类的有界代表与有限性}
\subsection*{来源与整理范围}
Andrew V. Sutherland, \emph{18.785 Number Theory I}, Fall 2021, Lecture 14: \emph{The Geometry of Numbers}, dated 10/27/2021. MIT OpenCourseWare. 本题目标为 Theorem 14.20 及 Corollary 14.22；收录 Proposition 14.16 至 Corollary 14.22 的完整证明链（印刷第5--8页），并摘列第4--5页的必要记号与测度规范。

实际访问的原文：\url{https://ocw.mit.edu/courses/18-785-number-theory-i-fall-2021/mit18_785f21_lec14.pdf}。
获取日期：2026-09-28。使用的是明确可复用的 OCW 2021 版，而非种子题单所列 2025 版；本文页码和编号均属于实际使用版本。

\subsection*{作者原命题与人类证明}
\paragraph{Notation from \S14.2 (equations (4)--(7)).}
Let $n=[K:\mathbb Q]=r+2s$. The embeddings identify $K_{\mathbb R}=K\otimes_{\mathbb Q}\mathbb R$ with the real subspace of $K_{\mathbb C}\simeq\mathbb C^n$ whose coordinates have the form
\[
(w_1,\ldots,w_r,x_1+iy_1,x_1-iy_1,\ldots,x_s+iy_s,x_s-iy_s).\tag{5}
\]
For $x,y\in K$, the inner product is
\[
\langle x,y\rangle=\sum_{\sigma}\sigma(x)\overline{\sigma(y)}.\tag{4}
\]
In these real coordinates, the normalized Haar measure is
\[
\mu(S)=2^s\mu_{\mathbb R^n}(S).\tag{6}
\]
For a lattice basis $e_1,\ldots,e_n$, with fundamental parallelepiped
$F(e_1,\ldots,e_n)=\{\sum t_i e_i:0\le t_i<1\}$, it satisfies
\[
\mu(F(e_1,\ldots,e_n))=\sqrt{|\det[\langle e_i,e_j\rangle]_{ij}|}.\tag{7}
\]

\begin{proposition}[14.16]
Let $K$ be a number field. Using the normalized Haar measure on $K_{\mathbb R}$ defined in (6),
\[
\covol(\mathcal O_K)=\sqrt{|D_K|}.
\]
\end{proposition}
\begin{proof}
Let $e_1,\ldots,e_n\in\mathcal O_K$ be a $\mathbb Z$-basis for $\mathcal O_K$, let $\Hom_{\mathbb Q}(K,\mathbb C)=\{\sigma_1,\ldots,\sigma_n\}$, and define $A:=[\sigma_i(e_j)]_{ij}\in\mathbb C^{n\times n}$. Then $D_K=\operatorname{disc}(e_1,\ldots,e_n)=(\det A)^2$, by Proposition 12.6.

Viewing $\mathcal O_K\hookrightarrow K_{\mathbb R}$ as a lattice in $K_{\mathbb R}$ with basis $e_1,\ldots,e_n$, we may use (7) to compute $\covol(\mathcal O_K)=\mu(F(e_1,\ldots,e_n))=\sqrt{|\det[\langle e_i,e_j\rangle]_{ij}|}$. Applying (4) yields
\[
\det[\langle e_i,e_j\rangle]_{ij}
=\det\left[\sum_k\sigma_k(e_i)\overline{\sigma_k(e_j)}\right]_{ij}
=\det(A^t\overline A)=(\det A)(\overline{\det A}).
\]
Noting that $\det A$ is the square root of an integer (hence either real or purely imaginary), we have $\covol(\mathcal O_K)^2=|(\det A)^2|=|D_K|$, and the proposition follows.
\end{proof}

Recall from Remark 6.13 that for number fields $K$ we view the absolute norm
\[
N:\mathcal I_{\mathcal O_K}\longrightarrow\mathcal I_{\mathbb Z},\qquad
I\longmapsto[\mathcal O_K:I]_{\mathbb Z}
\]
as having image in $\mathbb Q_{>0}$ by identifying $N(I)\in\mathcal I_{\mathbb Z}$ with a positive generator for $N(I)$ (note that $\mathbb Z$ is a PID). Recall that $[\mathcal O_K:I]_{\mathbb Z}$ is a module index of $\mathbb Z$-lattices in the $\mathbb Q$-vector space $K$ (see Definitions 6.1 and 6.5), and for ideals $I\subseteq\mathcal O_K$ this is just the positive integer $[\mathcal O_K:I]_{\mathbb Z}=[\mathcal O_K:I]$. When $I=(a)$ is a principal fractional ideal with $a\in K$, we may simply write $N(a):=N((a))=|N_{K/\mathbb Q}(a)|$.

\begin{corollary}[14.17]
Let $K$ be a number field and let $I$ be a nonzero fractional ideal of $\mathcal O_K$. Then
\[
\covol(I)=N(I)\sqrt{|D_K|}.
\]
\end{corollary}
\begin{proof}
Let $n=[K:\mathbb Q]$. Since $\covol(bI)=b^n\covol(I)$ and $N(bI)=b^nN(I)$ for any $b\in\mathbb Z_{>0}$, without loss of generality we may assume $I\subseteq\mathcal O_K$ (replace $I$ with a suitable $bI$ if not). Applying Propositions 14.11 and 14.16, we have
\[
\covol(I)=[\mathcal O_K:I]\covol(\mathcal O_K)=N(I)\covol(\mathcal O_K)=N(I)\sqrt{|D_K|}
\]
as claimed.
\end{proof}

\begin{theorem}[14.18; Minkowski bound]
Let $K$ be a number field of degree $n$ with $s$ complex places. Define the Minkowski constant $m_K$ for $K$ as the positive real number
\[
m_K:=\frac{n!}{n^n}\left(\frac4\pi\right)^s\sqrt{|D_K|}.
\]
For every nonzero fractional ideal $I$ of $\mathcal O_K$ there is a nonzero $a\in I$ for which
\[
N(a)\le m_KN(I).
\]
\end{theorem}
To prove this theorem we need the following lemma.

\begin{lemma}[14.19]
Let $K$ be a number field of degree $n$ with $r$ real and $s$ complex places. For each $t\in\mathbb R_{>0}$, the measure of the convex symmetric set
\[
S_t:=\left\{(z_\sigma)\in K_{\mathbb R}:\sum|z_\sigma|\le t\right\}\subseteq K_{\mathbb R}
\]
with respect to the normalized Haar measure $\mu$ on $K_{\mathbb R}$ is
\[
\mu(S_t)=2^r\pi^s\frac{t^n}{n!}.
\]
\end{lemma}
\begin{proof}
As in (5), we may uniquely write each $z=(z_\sigma)\in K_{\mathbb R}$ in the form
\[
(w_1,\ldots,w_r,x_1+iy_1,x_1-iy_1,\ldots,x_s+iy_s,x_s-iy_s)
\]
with $w_i,x_j,y_j\in\mathbb R$. We will have $\sum_\sigma|z_\sigma|\le t$ if and only if
\[
\sum_{i=1}^r|w_i|+\sum_{j=1}^s2\sqrt{|x_j|^2+|y_j|^2}\le t.\tag{8}
\]
We now compute the volume of this region in $\mathbb R^n$ by relating it to the volume of the simplex
\[
U_t:=\{(u_1,\ldots,u_n)\in\mathbb R_{\ge0}^n:u_1+\cdots+u_n\le t\}\subset\mathbb R^n,
\]
which is $\mu_{\mathbb R^n}(U_t)=t^n/n!$ (volume of the standard simplex in $\mathbb R^n$ scaled by a factor of $t$).

If we view all the $w_i,x_j,y_j$ as fixed except the last pair $(x_s,y_s)$, then $(x_s,y_s)$ ranges over a disk of some radius $d\in[0,t/2]$ determined by (8). If we replace $(x_s,y_s)$ with $(u_{n-1},u_n)$ ranging over the triangular region bounded by $u_{n-1}+u_n\le2d$ and $u_{n-1},u_n\ge0$, we need to incorporate a factor of $\pi/2$ to account for the difference between $(2d)^2/2=2d^2$ and $\pi d^2$; repeat this $s$ times. Similarly, if we hold everything but $w_r$ fixed and replace $w_r$ ranging over $[-d,d]$ for some $d\in[0,t]$ with $u_r$ ranging over $[0,d]$, we need to incorporate a factor of $2$ to account for this change of variable; repeat $r$ times. We then have
\[
\mu(S_t)=2^s\mu_{\mathbb R^n}(S_t)=2^s\left(\frac\pi2\right)^s2^r\mu_{\mathbb R^n}(U)=2^r\pi^s\frac{t^n}{n!}.
\]
\end{proof}

\begin{proof}[Proof of Theorem 14.18]
Let $I$ be a nonzero fractional ideal of $\mathcal O_K$. By Theorem 14.13, if we choose $t$ so that $\mu(S_t)>2^n\covol(I)$, then $S_t$ will contain a nonzero $a\in I$. By Lemma 14.19 and Corollary 14.17, it suffices to choose $t$ so that
\[
\left(\frac tn\right)^n
=\frac{n!\mu(S_t)}{n^n2^r\pi^s}
>\frac{n!2^n}{n^n2^r\pi^s}\covol(I)
=\frac{n!}{n^n}\left(\frac4\pi\right)^s\sqrt{|D_K|}N(I)
=m_KN(I).
\]
Let us now pick $t$ so that $(t/n)^n>m_KN(I)$. Then $S_t$ contains $a\in I$ with $\sum_\sigma|\sigma(a)|\le t$. Recalling that the geometric mean is bounded above by the arithmetic mean, we then have
\[
N(a)=\left(N(a)^{1/n}\right)^n
=\left(\prod_\sigma|\sigma(a)|^{1/n}\right)^n
\le\left(\frac1n\sum_\sigma|\sigma(a)|\right)^n
\le\left(\frac tn\right)^n.
\]
Taking the limit as $(t/n)^n\to m_KN(I)$ from above yields $N(a)\le m_KN(I)$.
\end{proof}

\paragraph{14.5 Finiteness of the class group.}
Recall that the ideal class group $\operatorname{cl}\mathcal O_K$ is the quotient of the ideal group $\mathcal I_K$ of $\mathcal O_K$ by its subgroup of principal fractional ideals. We now use the Minkowski bound to prove that every ideal class $[I]\in\operatorname{cl}\mathcal O_K$ can be represented by an ideal $I\subseteq\mathcal O_K$ of small norm. It will then follow that the ideal class group is finite.

\begin{theorem}[14.20]
Let $K$ be a number field. Every ideal class in $\operatorname{cl}\mathcal O_K$ contains an ideal $I\subseteq\mathcal O_K$ of absolute norm $N(I)\le m_K$, where $m_K$ is the Minkowski constant for $K$.
\end{theorem}
\begin{proof}
Let $[J]$ be an ideal class of $\mathcal O_K$ represented by the nonzero fractional ideal $J$. By Theorem 14.18, the fractional ideal $J^{-1}$ contains a nonzero element $a$ for which
\[
N(a)\le m_KN(J^{-1})=m_KN(J)^{-1},
\]
and therefore $N(aJ)=N(a)N(J)\le m_K$. We have $a\in J^{-1}$, thus $aJ\subseteq J^{-1}J=\mathcal O_K$, so $I=aJ$ is an $\mathcal O_K$-ideal in the ideal class $[J]$ with $N(I)\le m_K$ as desired.
\end{proof}

\begin{lemma}[14.21]
Let $K$ be a number field of degree $n$ and let $M\in\mathbb R_{>0}$. The number of $\mathcal O_K$-ideals of norm $N(I)\le M$ is bounded by $(nM)^{\log_2 M}$ (and in particular, finite).
\end{lemma}
\begin{proof}
Let $I$ be an ideal of absolute norm $N(I)\le M$ and let $I=\mathfrak p_1\cdots\mathfrak p_k$ be its factorization into (not necessarily distinct) prime ideals. We have $N(\mathfrak p_i)\ge2$ for each $\mathfrak p_i$ so $k\le\log_2 M$. There are less than $M$ primes $p\le M$ and at most $n$ primes $\mathfrak p$ of $\mathcal O_K$ above each $p\le M$. It follows that the there are less than $(nM)^{\log_2 M}$ $\mathcal O_K$-ideals with norm $N(I)\le M$.
\end{proof}

\begin{corollary}[14.22]
Let $K$ be a number field. The ideal class group of $\mathcal O_K$ is finite.
\end{corollary}
\begin{proof}
By Theorem 14.20, each ideal class is represented by an ideal of norm at most $m_K$, and by Lemma 14.21, the number of such ideals is finite.
\end{proof}


\subsection*{整理说明（非作者正文）}
本题对应种子题单第6题。$s$ 即种子题单的 $r_2$，$D_K$ 即 $\Delta_K$；作者原记号不改。前置 ``Notation'' 段是整理者从原文 (4)--(7) 摘编的记号说明，不是新增人类证明；随后命题、证明和过渡段按原顺序转录。

标准先修为 Minkowski 凸体定理（原文 Theorem 14.13）、格指数与余体积关系（Proposition 14.11）、判别式的嵌入矩阵公式（Proposition 12.6）、Dedekind 环理想分解及范数性质。关键常数的体积推导、逆理想构造、有限性论证均收在本题中，没有作为先修跳过。

仅统一了数学排版及断行；原文 Lemma 14.19 最后一式写 $U$，保留该写法（同段先定义 $U_t$）；Lemma 14.21 中的 ``the there'' 亦照录，不对原文另作数学审稿。仅保留与本题对应的结论，不收入原文将证明留作习题的 Theorem 14.23。PDF 原页已实际打开，并逐页查看第5--8页图像；原 PDF 的本地下载未成功，包内仅保存本次助手转录的摘录，不将其冒称原 PDF。

\subsection*{可修改的简短标注（非作者正文）}
\textbf{$c$：}数域理想类与有界代表。

\textbf{$a$：}数域的实嵌入；判别式；格与余体积；Minkowski 凸体定理；凸体体积；算术--几何平均不等式；逆分式理想；素理想分解。

\textbf{cross\_theory\_status = yes。}原文 Proposition 14.16、Corollary 14.17 把判别式和理想范数转为格余体积；Theorem 14.18 再将凸体体积控制转回小范数元素，Theorem 14.20 用它取得每个理想类的有界代表。位置：印刷第5--8页。

这些标签是可修改的非穷尽初稿，不作正式跨理论认证。难度选取依据是本题需要连通嵌入、体积规范、常数计算与理想类代表构造，而非只直接引用 Minkowski 界。
\subsection*{许可与修改记录}
材料来自 MIT OpenCourseWare；原作者与课程见上文。按该站所示 \href{https://creativecommons.org/licenses/by-nc-sa/4.0/}{CC BY-NC-SA 4.0} 许可复用。本转录及其附加整理采用同一许可；不得据此声称作者或 MIT 认可本次整理。2026-09-28：ChatGPT 转录为 TeX，加入题号、来源、整理说明和初稿标注；数学正文保持作者语言及顺序。

\end{document}
